% ─────────────────────────────────────────────────────────────────────────────
\section{Operational Interpretation and Experimental Limitations}
\label{sec:results-limitations}
% ─────────────────────────────────────────────────────────────────────────────

The results demonstrate that the operational robustness of \gls{ai}-based
systems for forensic image classification and triage cannot be inferred from
clean accuracy alone. A complete interpretation must jointly consider nominal
performance, error direction, behavior on the separate \gls{ood} branch,
sensitivity to adversarial and anti-forensic manipulations, configuration and
mapping dependence, and the extent to which each output remains traceable and
subject to human verification.

False negatives and false positives have different consequences within a
forensic workflow. A false negative on the \texttt{weapon} class may prevent a
potentially relevant image from being prioritized for review. A false positive
does not conceal target material but may increase analyst workload, dilute
prioritization, and reduce the practical selectivity of the automated process.

A \texttt{weapon} assignment on an \gls{ood} input is not a binary false
positive because the input has no supervised
\texttt{weapon}/\texttt{non\_weapon} reference assignment. Nevertheless, such
an output may produce similar operational effects by routing unrelated,
ambiguous, degraded, synthetic, or otherwise out-of-task material toward the
firearm-oriented review category.

Aggregate accuracy therefore provides only a partial description of the
observed operating profile. Similar accuracy values may conceal substantially
different combinations of missed target images, additional review workload,
and changed individual decisions. Conversely, unchanged or slightly increased
aggregate accuracy does not establish sample-level invariance, because newly
introduced errors may be offset by previously incorrect decisions that become
correct.

The parallel evaluation also shows that nominally strong systems may fail under
different conditions. EfficientNet-B0 achieves the highest clean accuracy among
the proxy models but undergoes severe direct-target degradation under
\gls{sigmazero} and substantial degradation under \gls{fgsm}. ResNet18 and the
\gls{clip}-based proxy exhibit more limited effects from the same
EfficientNet-B0-oriented artifacts, but these results measure transfer behavior
rather than direct robustness against attacks optimized for those
architectures.

The selected black-box software configurations do not reproduce the
EfficientNet-B0 direct-target collapse. Their normalized outputs nevertheless
reveal other operational limitations, including missed detections,
false-positive workload, configuration sensitivity, non-zero \gls{ood} weapon
assignment rates, and transformation-specific changes in the balance between
the two binary error directions.

These differences reinforce the need to evaluate automated image-analysis
systems across multiple operating conditions rather than relying on a single
clean benchmark. They also show that robustness must be understood as a
property of a specific system version, configuration, task definition, input
population, output mapping, and threat model.

The following limitations delimit the scope of the experimental findings.

\begin{itemize}

\item \textbf{Dataset size and intended scope.}
The frozen dataset contains 1\,500 source images: 500 \texttt{weapon} images,
500 \texttt{non\_weapon} images, and 500 images reserved for the separate
\gls{ood} branch. This population supports a controlled operational robustness
study but does not represent the complete diversity of firearm-related,
non-firearm, ambiguous, degraded, synthetic, and out-of-task imagery that may
occur in operational investigations.

\item \textbf{Artificial binary balance.}
The binary subset contains equal numbers of \texttt{weapon} and
\texttt{non\_weapon} images. This balance supports controlled comparison of the
two error directions, but it must not be interpreted as an estimate of class
prevalence in real forensic collections. Operational accuracy, precision, and
review workload may change under substantially different class distributions.

\item \textbf{Human semantic review.}
The final operational assignments derive from a documented
\textit{human-in-the-loop} selection procedure conducted by a single reviewer.
No independent second annotation or inter-rater agreement measurement was
included. The labels are therefore controlled and traceable but remain
dependent on the task definitions and judgment applied during the review
process.

\item \textbf{Source-distribution imbalance.}
The source population is heterogeneous but not uniformly distributed across
the final classes. In particular, the \texttt{non\_weapon} branch is
predominantly derived from the \texttt{05\_deepweb} source group.
Source-specific visual regularities may therefore influence some observed
classification behavior despite source-aware fold construction.

\item \textbf{Heterogeneous \gls{ood} branch.}
The \gls{ood} population intentionally combines borderline, ambiguous,
degraded, synthetic, non-firearm weapon, and semantically out-of-task content.
It is not a homogeneous class and was not used as a third supervised training
label. The reported weapon assignment rates summarize closed-set behavior over
this frozen mixture and must not be interpreted as performance on a
well-defined universal \gls{ood} distribution.

\item \textbf{Derived-sample dependence.}
The 5\,000 adversarial and 5\,000 anti-forensic artifacts are derived from the
same 1\,000 clean binary source images. The 11\,000-item binary evaluation
population therefore does not represent 11\,000 independent underlying cases.
Bundle-level metrics characterize behavior across the frozen set of
source--condition combinations and must not be interpreted as independent
population-level observations.

\item \textbf{Finite perturbation suite.}
The adversarial and anti-forensic suites include a finite set of conditions,
with one frozen parameterization for each attack or transformation. Different
perturbation budgets, optimization parameters, image encodings, transformation
severities, software implementations, or sequential combinations may produce
different robustness profiles.

\item \textbf{Restricted direct-target adversarial evaluation.}
The four model-dependent adversarial attacks were generated against
fold-specific EfficientNet-B0 checkpoints. Their evaluation on EfficientNet-B0
constitutes direct-target assessment under the documented access assumptions.
Their evaluation on ResNet18, the \gls{clip}-based proxy, and the selected
black-box software configurations measures only empirical transfer of the
frozen artifacts.

The experiment therefore does not characterize attacks optimized independently
against ResNet18, the \gls{clip}-based proxy, or any commercial software system.
Limited transfer must not be interpreted as general direct-attack robustness.

\item \textbf{Proxy-model role.}
EfficientNet-B0, ResNet18, and the \gls{clip}-based binary-head model are
transparent experimental references. They permit controlled access to
predictions, fold-specific checkpoints, gradients, and maximum predicted-class
probabilities. They must not be interpreted as reconstructions, surrogates with
demonstrated fidelity, or architectural approximations of the internal models
used by the selected proprietary software systems.

\item \textbf{Score interpretation and calibration.}
The maximum predicted-class probabilities produced by the proxy models are
model-specific outputs and are not treated as calibrated estimates of
correctness or forensic reliability. Their numerical values must not be
compared across architectures as though they shared a common uncertainty scale.

The selected black-box software systems do not expose a homogeneous numerical
score suitable for an equivalent cross-system analysis. Exported percentages,
tags, retrieval membership, classifications, and bookmarks have different
semantics and cannot be converted into directly comparable calibrated
probabilities.

\item \textbf{Different \gls{ood} prediction units.}
Each proxy-model \gls{ood} rate aggregates 2\,500 fold-specific responses,
corresponding to five checkpoint evaluations for each of 500 unique images.
Each black-box rate uses one normalized image-level decision for each of the 500
images.

The proportions support a parallel descriptive comparison of operational
assignment tendencies, but they do not constitute identical statistical
observations or directly interchangeable denominators.

\item \textbf{Black-box version and configuration dependence.}
The black-box findings are bounded to the evaluated software versions, enabled
processing functions, imported population, export procedures, and explicit
normalization rules. Updates to internal models, preprocessing, category
definitions, retrieval mechanisms, or default settings may change the
observable results.

\item \textbf{Operational normalization of heterogeneous outputs.}
The common \texttt{weapon}/\texttt{non\_weapon} representation is an explicit
operational abstraction constructed from observable software outputs. It does
not establish that the selected systems implement equivalent internal binary
tasks.

The Magnet.AI results depend on the \texttt{Possible weapons} tag. The
Griffeye/T3K CORE results depend on the conservative use of
\texttt{CORE/Violence/Firearm} as the primary positive bookmark. The Cellebrite
results depend on the union of the exported \texttt{Armi}, \texttt{Pistola},
and \texttt{Fucile} labels. The Excire results depend on the eight
firearm-oriented semantic queries, union-based retrieval mapping, and selected
\texttt{Distance limit}.

Alternative defensible mappings would define different operational tasks and
could alter the observed balance between missed detections and false-positive
workload.

\item \textbf{Excire configuration sensitivity.}
Excire Foto 2025 was evaluated as a semantic image-retrieval system rather than
as a native binary forensic classifier. The \texttt{D20}, \texttt{D50}, and
\texttt{D80} settings produce materially different recall, \gls{fpr}, and
\gls{ood} assignment profiles. The intermediate \texttt{D50} configuration
facilitates cross-system comparison but is not established as a universally
optimal operating point.

\item \textbf{Absence of internal black-box diagnostics.}
The experiment cannot identify the proprietary architectures, training data,
feature representations, preprocessing stages, thresholds, or internal
decision mechanisms responsible for the observed black-box behavior. The
results characterize normalized observable outputs rather than internal causal
mechanisms.

\item \textbf{Qualitative \gls{xai} scope.}
The retained \gls{xai} analysis is restricted to five purposively selected
EfficientNet-B0 cases. The cases illustrate different operating conditions and
failure modes but do not estimate their frequency in the complete population.

The historical \gls{ig} assets were generated using 32 integration steps and a
zero baseline in the preprocessed input space. Numerical convergence was not
established for the retained figures, and the cases remain pending adaptive
regeneration under the hardened validation procedure. The maps are therefore
used only as qualitative diagnostic representations.

They must not be interpreted as semantic segmentations, causal explanations,
proofs of feature correctness, or explanations of the selected black-box
software systems.

\item \textbf{Embedded-metadata sensitivity.}
The metadata-preserving blind bundle concealed reference information in file
names, paths, and directory structure but did not remove metadata embedded
within the image bytes. The audit identified 15 files containing predefined
semantic or experiment-related sensitive-term matches.

The leave-out analysis indicates that excluding the affected clean and
\gls{ood} cases produces only limited aggregate metric variation. It does not,
however, establish whether any evaluated black-box system accessed or used the
metadata. No metadata-neutralized paired bundle was included, and causal
separation of visual and metadata contributions is therefore outside the scope
of the experiment.

\item \textbf{Controlled-data and public-artifact boundary.}
Complete raw image corpora, generated image trees, proprietary export contents,
and commercial prediction-level outputs cannot all be redistributed through
the public repository because of source-specific, legal, ethical,
institutional, and proprietary restrictions.

The repository preserves scripts, schemas, manifests, integrity records,
normalization summaries, metric tables, selected reporting assets, and
documentation required to audit the experimental process. This supports
procedural reproducibility and artifact-level traceability, but complete public
re-execution still requires authorized restoration of the controlled datasets,
checkpoints, and proprietary software outputs.

\end{itemize}

These limitations define the conditions under which the experimental
conclusions are supported. The chapter provides a protocol-specific assessment
of operational robustness rather than:

\begin{itemize}
\item certification of any evaluated software product;
\item a universal ranking of models or commercial systems;
\item a complete adversarial-security assessment;
\item validation of a production-grade firearm detector;
\item evidentiary validation of automated classifications;
\item authorization for autonomous operational deployment.
\end{itemize}

The practical implication is that automated image-analysis systems should be
used as traceable triage support within a \textit{human-in-the-loop} forensic
workflow. Their outputs should remain associated with the source item, system
or model version, checkpoint or software configuration, observable mapping,
evaluated condition, and relevant limitations.

Missed detections, false-positive workload, \gls{ood} assignments,
configuration-dependent behavior, and failures accompanied by high maximum
predicted-class probabilities must remain visible to the analyst rather than
being concealed behind a single aggregate metric.

Operational deployment would require additional task-specific validation on
representative local data, documented change management, periodic reevaluation
after software or model updates, and explicit procedures for reviewing and
overriding automated outputs.

The final classification or evidentiary interpretation must remain the
responsibility of the qualified human analyst. The evaluated systems may assist
with prioritization and large-scale review, but their outputs do not replace
contextual assessment, corroboration, chain-of-custody requirements, or expert
judgment.
